% ---------------------------------------------------------------------
\section{扭度、规范不变算符与欧氏算符乘积展开}\label{sec:dis-ope}

当 $q$ 为大的欧氏动量时，$\ii\int\dd^4z\,\ee^{\ii q\cdot z}T\,j^\mu(z)j^\nu(0)$ 由 $z\to0$ 的区域主导，两流之间的硬传播子收缩为一点，乘积展开为定域算符之和，系数只依赖于 $q$，矩阵元只依赖于靶。对前向矩阵元，自旋 $n$ 的算符给出 $\langle P|O^{\mu_1\cdots\mu_n}|P\rangle\propto P^{\mu_1}\cdots P^{\mu_n}$，与系数中的 $q_{\mu_1}\cdots q_{\mu_n}/(Q^2)^n$ 缩并后正是 $(P\cdot q/Q^2)^n\propto\omega^n$：OPE 的第 $n$ 项给出 Taylor 系数，量纲分析给出它的幂次 $(Q^2)^{-(d-n-2)/2}$，故领头项来自扭度 $d-n=2$ 的算符。矩只需要这一\emph{短距离}展开在单个 Taylor 系数上的形式；\cref{fig:ope-schematic}所画的\emph{光锥}展开（每一扭度包含所有自旋）是 $x$ 空间的结构，留到\cref{prop:dis-light-cone}讨论。

%%INSERT fig:ope-schematic

\begin{definition}[扭度与扭度 2 算符]\label{def:dis-twist}
(a) 对每个味 $q$ 与 $n\ge1$，以及（胶子）$n\ge2$，令
\begin{equation}\label{eq:dis-twist2-ops}
O_q^{\mu_1\cdots\mu_n}=\mathcal S\,\bar\psi_q\gamma^{\mu_1}\ii D^{\mu_2}\cdots\ii D^{\mu_n}\psi_q-\text{迹},\qquad
O_g^{\mu_1\cdots\mu_n}=\mathcal S\,G^{\mu_1\alpha}\ii D^{\mu_2}\cdots\ii D^{\mu_{n-1}}G_\alpha{}^{\mu_n}-\text{迹},
\end{equation}
$\mathcal S$ 为对 $\mu_1,\dots,\mu_n$ 的对称化，“$-$迹”使张量无迹，胶子算符中的协变导数取伴随表示。$O_g$ 的符号由 $O_g^{+\cdots+}=\sum_jG^{+j}(\ii D^+)^{n-2}G^{+j}$ 固定，于是对在壳、物理极化平均的胶子 $\langle g(k)|O_g^{+\cdots+}|g(k)\rangle=2(k^+)^n$。% VF-CONV-15
\par
(b) 约化矩阵元 $A^{(n)}_i$（$i=q,g$）由 $\langle P|O_i^{\mu_1\cdots\mu_n}|P\rangle=2A^{(n)}_i\big(P^{\mu_1}\cdots P^{\mu_n}-\text{迹}\big)$ 定义（对靶自旋平均）。自由夸克靶的树图值为 $A^{(n)}_q=1$。
\par
(c) $O_q^{+\cdots+}=\bar\psi\gamma^+(\ii D^+)^{n-1}\psi$ 的零胶子与单胶子顶点为
\begin{equation}\label{eq:dis-Oq-vertex}
\gamma^+(k^+)^{n-1},\qquad g\,t^a\,n_B^\mu\,\gamma^+\sum_{j=0}^{n-2}(k_2^+)^j(k_1^+)^{n-2-j},
\end{equation}
$k$ 为夸克动量，$k_1$、$k_2=k_1+\ell$ 为流入与流出的夸克动量，$\ell$ 为流入的胶子动量（$\ii D^+=\ii\partial^++gA^+$）。% VF-PDF-12
\par
(d) 对称无迹、秩为 $n$、典型量纲为 $d$ 的算符，其扭度为 $\tau=d-n$；\cref{eq:dis-twist2-ops}的算符 $d=n+2$，扭度为 2。场分量的\emph{共线扭度}定义为 $t=d-h$，其中推动权 $h$ 是沿 $+z$ 推动 $v^\pm\to\ee^{\pm\eta}v^\pm$ 时该分量所乘的 $\ee^{h\eta}$ 的指数：每个上指标 $+$ 记 $+1$，$-$ 记 $-1$，横向记 $0$；好分量 $\psi_+=\mathcal P_+\psi$ 记 $+\frac12$，坏分量 $\psi_-=\mathcal P_-\psi$ 记 $-\frac12$（$\mathcal P_\pm=\frac12\gamma^\mp\gamma^\pm$）。$t$ 对乘积可加，且对称无迹算符的 $+\cdots+$ 分量的共线扭度等于其扭度 $\tau$。
\uses{def:prelim-qcd}
\end{definition}

\begin{lemma}[扭度 2 规范不变算符的分类 \tagR]\label{lem:dis-twist-classification}
（$R_{\mathrm{def}}$：由算符定义直接推出。）构件的共线扭度为
\begin{equation}\label{eq:dis-twist-table}
\begin{array}{c|cccccccc}
 & \psi_+,\ \bar\psi_+ & \psi_- & D^+ & D^j & D^- & F^{+j} & F^{+-},\ F^{ij} & F^{-j}\\\hline
t & 1 & 2 & 0 & 1 & 2 & 1 & 2 & 3
\end{array}
\end{equation}
（$\bar\psi_+\equiv\bar\psi\mathcal P_-$）。在全导数与运动方程算符之外，具有非零的自旋平均前向矩阵元的扭度 2 对称无迹规范不变算符由 $O^{(n)}_q$（每个味一个）与 $O^{(n)}_g$（$n$ 为偶数）张成。（本引理是算符的代数陈述。）
\uses{def:dis-twist,lem:prelim-boost}
\end{lemma}

\begin{proof}
完整证明见附录~M（\cref{sec:appM-brst}），要点如下。(1) 由\cref{lem:prelim-boost}，$\psi_\pm$ 乘以 $\ee^{\pm\eta/2}$，矢量分量 $v^\pm$ 乘以 $\ee^{\pm\eta}$；与 $d(\psi)=\frac32$、$d(D)=1$、$d(F)=2$ 合起来即得表\cref{eq:dis-twist-table}。(2) 规范不变的色单态单项式至少含两个基本场（$\psi$、$\bar\psi$ 或 $F$），每个基本场 $t\ge1$，而导数 $t\ge0$ 且只有 $D^+$ 的 $t=0$；故总扭度为 2 时恰有两个 $t=1$ 的场（$\psi_+$、$\bar\psi_+$ 或 $F^{+j}$），导数全是 $D^+$。(3) 夸克双线性型 $\bar\psi\mathcal P_-\Gamma\mathcal P_+\psi$ 只对 $\Gamma\in\{\gamma^+,\gamma^+\gamma_5,\gamma^+\gamma^j\}$ 不为零（$\mathcal P_-\Gamma\mathcal P_+$ 张成四维空间）；$\gamma^+\gamma_5$ 给出赝张量，宇称使其自旋平均的非极化前向矩阵元为零；$\gamma^+\gamma^j$ 带一个不能被缩并的横向指标，属于混合对称表示（横度），且前向自旋平均的矩阵元中没有可用的横向矢量。味不对角的双线性型在味本征态中的前向矩阵元为零。(4) 胶子 $F^{+j}(D^+)^{n-2}F^{+k}$ 中 $\delta^{jk}$ 给出 $O_g$，$\epsilon^{jk}$ 给出赝张量，$jk$ 的对称无迹部分需要横向张量，二者的非极化前向矩阵元都为零。(5) 把 $D^+$ 从一个场移到另一个场只差全导数，前向矩阵元为零；$n$ 为奇数时胶子算符（模全导数）是反对称的，即全导数。
\end{proof}

\begin{lemma}[扭度奇偶性与手征 \tagR]\label{lem:dis-twist-chirality}
（$R_{\mathrm{def}}$：由算符定义直接推出。）设 $O$ 是 $\psi$、$\bar\psi$、$A$ 及其导数的玻色型、不含鬼场的规范不变定域多项式，在 $\psi\to\gamma_5\psi$、$\bar\psi\to-\bar\psi\gamma_5$ 下的本征值为 $\chi=\pm1$（手征），属于秩为 $j$ 的对称无迹表示，典型量纲为 $d$。则 $(-1)^{d-j}=\chi$：手征偶的算符扭度为偶数，扭度为奇数的算符是手征奇的。\emph{注意：}对混合对称的表示这一结论不成立——横度算符 $\bar\psi\,\ii\sigma^{+j}\psi$ 的扭度为 2，却是手征奇的。
\uses{def:dis-twist,lem:dis-twist-classification}
\end{lemma}

\begin{proof}
完整证明见附录~M（\cref{sec:appM-brst}）。由表\cref{eq:dis-twist-table}，玻色型构件（$D^\pm,D^j,F^{\mu\nu}$）的共线扭度与其横向指标数 $n_T$ 同奇偶：$t\equiv n_T\pmod 2$。对夸克双线性型，手征偶（$\Gamma$ 含奇数个 $\gamma$）时 $t\equiv n_T$，手征奇时 $t\equiv n_T+1$（例如 $\bar\psi\gamma^+\psi$：$t=2$，$n_T=0$；$\bar\psi\psi=\bar\psi_+\psi_-+\bar\psi_-\psi_+$：$t=3$，$n_T=0$）。在对称无迹算符的 $+\cdots+$ 分量中，自由指标都是 $+$，横向指标只能两两缩并（$\delta^{jk}$ 或 $\epsilon^{jk}$），故 $\sum n_T$ 为偶数，总扭度的奇偶等于手征奇双线性型数目的奇偶，即 $(-1)^{\tau}=\chi$，而 $\tau=d-j$。混合对称的算符有横向的自由指标，$\sum n_T$ 可以为奇数，论证失效。
\end{proof}

对称无迹算符的系数张量由协变性与流守恒确定：含 $g_{\mu_i\mu_j}$ 的项与无迹算符缩并为零；横向化以后只剩两个结构，记
\begin{equation}\label{eq:dis-ope-tensor}
T^{\mu\nu}\simeq\sum_{n=2,4,\dots}\sum_i\Big[\Big(-g^{\mu\nu}+\frac{q^\mu q^\nu}{q^2}\Big)\frac{2^nq_{\mu_1}\cdots q_{\mu_n}}{(Q^2)^n}\,C^{(n)}_{1,i}
+\Big(g^\mu{}_{\mu_1}-\frac{q^\mu q_{\mu_1}}{q^2}\Big)\Big(g^\nu{}_{\mu_2}-\frac{q^\nu q_{\mu_2}}{q^2}\Big)\frac{2^nq_{\mu_3}\cdots q_{\mu_n}}{(Q^2)^{n-1}}\,C^{(n)}_{2,i}\Big]O_i^{\mu_1\cdots\mu_n}.
\end{equation}
\emph{收缩规则：}以\cref{def:dis-twist}(b) 的矩阵元代入并略去迹项（无质量靶时迹项为零），第一项给出 $-\tilde g^{\mu\nu}(2P\cdot q/Q^2)^n\,2A^{(n)}_iC^{(n)}_{1,i}$，第二项给出 $\tilde P^\mu\tilde P^\nu\,2^n(P\cdot q)^{n-2}(Q^2)^{1-n}\,2A^{(n)}_iC^{(n)}_{2,i}$；与\cref{eq:dis-T}比较，
\begin{equation}\label{eq:dis-contraction}
T_1\ni2\,\omega^nC^{(n)}_{1,i}A^{(n)}_i,\qquad T_2\ni4\,\omega^{n-1}C^{(n)}_{2,i}A^{(n)}_i
\end{equation}
（\texttt{s2\_compton.py} 对 $n=2,4,6$ 核对）。% VF-DIS-13
下面的引理是这一结构最简单的实例。

\begin{lemma}[树图 OPE \tagR]\label{lem:dis-tree-ope}
（$R_{\mathrm{PT}}$：在重整化微扰论中对红外正规化的部分子态严格。）在树图阶、无质量夸克态之间，$T^{\mu\nu}$ 由\cref{eq:dis-ope-tensor}给出（无余项），其中对偶数 $n$ 有 $C^{(n)}_{1,q}=C^{(n)}_{2,q}=e_q^2$；奇数 $n$ 的项在 $s$ 道与 $u$ 道图之间相消（交叉对称性）。（本引理只涉及部分子态。）
\uses{def:dis-twist,def:prelim-feynman-rules,def:dis-compton,ex:dis-free-quark}
\end{lemma}

\begin{proof}
由\cref{ex:dis-free-quark}，$T_2=4e_q^2\omega/(1-\omega^2)=\sum_{n=2,4,\dots}4e_q^2\omega^{n-1}$，$T_1=2e_q^2\omega^2/(1-\omega^2)=\sum_{n=2,4,\dots}2e_q^2\omega^n$（$|\omega|<1$）。自由夸克的树图矩阵元为 $\langle p|O^{(n)}_q|p\rangle=\bar u\gamma^{\mu_1}u\,p^{\mu_2}\cdots p^{\mu_n}$ 的对称化 $=2p^{\mu_1}\cdots p^{\mu_n}$（$p^2=0$，无迹），即 $A^{(n)}_q=1$。由\cref{eq:dis-contraction}，逐项比较 $\omega$ 的幂次得 $C^{(n)}_{1,q}=C^{(n)}_{2,q}=e_q^2$（$n$ 偶）。由于两个不变振幅完全确定横向张量，这就是张量层次的等式。$s$ 道图单独给出 $T_2^{(s)}=2e_q^2/(1-\omega)=2e_q^2\sum_{k\ge0}\omega^k$，含 $\omega$ 的一切幂次；$u$ 道图是 $T_2^{(u)}(\omega)=-T_2^{(s)}(-\omega)$，二者相加恰好消去 $\omega^{n-1}$ 中 $n$ 为奇数的项。等价地，在 $s$ 道中把传播子 $\frac{1}{(p+q)^2}=\frac{-1}{Q^2(1-\omega)}$ 按 $\omega$ 展开，第 $k$ 项正是算符 $\bar\psi\gamma^{\{\mu_1}\ii\partial^{\mu_2}\cdots\ii\partial^{\mu_{k+1}\}}\psi$ 的树图矩阵元乘以 $q_{\mu}$ 的乘积（\cref{def:dis-twist}(c) 的零胶子顶点）。
\end{proof}
